%%%%%
%
\section{Transformation d'un temps GPS en temps UTC (version automatique synchrone)}

\itemstyle{Op\'erateur~:} \operfmt{GPS\_to\_UTC\_auto}
\index{GPS\_to\_UTC\_auto}
\label{OP-GPS_to_UTC_auto}

%%%%%OBJET%%%%
\itemstyle{Objet~:}
Le temps GPS est compt� en secondes depuis le d�but de
la semaine, dans le syst�me GPS (d�cal� d'une valeur
constante avec le TAI et d'une valeur variable avec l'UTC).
 Cet op�rateur
convertit des temps exprim�s dans le syst�me GPS en temps UTC,
compt� en secondes depuis 0h00.

Cet op�rateur est dit automatique, car il utilise 
la valeur du d�calage GPS-UTC fournie par
le r�cepteur.

%%%%%ARGUMENTS%%%%
\bigskip
%%%%%ENTREES%%%%
\begin{tabarg}
\multicolumn{3}{l}{\titretab {Entr\'ees}}\\
\hline
$T_{GPS}$ & Vecteur, \texttt{float} & Temps dans le syst�me GPS (second)\\
$Decal$ & Vecteur, \texttt{float} & Information sur le d�calage GPS/UTC tel
que fournie par le syst�me Trimble/Bancom utilis�s sur le Merlin, l'Aztec, l'ATR et le Falcon (second)\\
\end{tabarg}

\medskip
%%%%%SORTIES%%%%
\begin{tabarg}
\multicolumn{3}{l}{\titretab {Sorties}}\\
\hline
$T_{UTC}$ & Vecteur, \texttt{double} & $T_{GPS}$  exprim� dans le syst�me UTC (second)\\
\end{tabarg}


%%%%%FORMULE%%%%
\itemstyle{Formule}

%% GPS_to_UTC_auto
  Recherche des valeurs non-nulles de $Decal$, et
v�rification qu'elles sont identiques. Si c'est le
cas, on d�finit $D$ comme cette valeur.

  Sur le premier $T_{GPS}$ non-nul, calcul d'un offset : 
\[
offset = 86400 . [(T_{GPS}-D) / 86400]
\]
Puis :
\[
   T_{UTC} = T_{GPS} - D - offset
\]



%%%%%COMMENTAIRE%%%%
\itemstyle{Commentaire~:}
Si un vecteur des diff�rences GPS-UTC n'est pas disponible ou 
inutilisable, il faut utiliser l'op�rateur 
\hyperref[OP-GPS_to_UTC]{\operfmt{GPS\_to\_UTC}} avec un argument explicite

Si le vecteur des diff�rences GPS-UTC n'est pas synchrone avec
le vecteur des temps, il faut utiliser l'op�rateur 
\hyperref[OP-GPS_to_UTC_irreg]{\operfmt{GPS\_to\_UTC\_irreg}}.

%%%%%AUTEUR%%%%
\itemstyle{Auteur~:}
CNRM/GMEI/TRAMM

%%%%%REFERENCE%%%%
\itemstyle{R\'ef\'erences~:}
\textit{GPS, theory and practice}, par B. Hofmann-Wellenhof, H. Lichtenegger, 
and J. Collins \cite{Hofmann-GPS}
